\subsection{The actual terminal transfer at every head}
\label{actual-terminal-transfer}

The compact density theorem now allows a regular radial identity to
determine every orbital vector.  To use it, one must construct a compact
coefficient having the required terminal radial image.  The following
lift keeps its scalar normalization explicit.

\begin{lemma}\label{terminal-transfer-upper-cut}
For every $r\ge0$ there is a nonnegative smooth function $b_U$ on
$\operatorname{GL}_r(\mathbb C)$ such that
\[
       \int_{N_r}b_U(Zu)\,du=1\qquad(Z\in\operatorname{GL}_r(\mathbb C)),
\]
where $du$ is planar Lebesgue measure in strict upper entries.
Moreover $b_U$ cuts every compact subset modulo right $N_r$ to a
compact subset of $\operatorname{GL}_r(\mathbb C)$, also with any
additional parameter in an arbitrary topological space.
\end{lemma}

\begin{proof}
Use the unique smooth QR decomposition $Z=QAU$, where $Q$ is
unitary, $A$ is positive diagonal, and $U\in N_r$.  It follows
directly by Gram--Schmidt with positive square roots.  Right
multiplication by $u\in N_r$ leaves $Q,A$ fixed and replaces
$U$ by $Uu$.  Choose $b\in C_c^\infty(N_r)$, $b\ge0$,
$\int b=1$, and put $b_U(Z)=b(U(Z))$.  Left invariance of the
upper-entry measure proves the integral formula.

For a compact $K\subset\operatorname{GL}_r(\mathbb C)\times X$,
the image of $K\times\operatorname{supp}b$ under
\[
       ((Z_0,x),u_0)\longmapsto(Q(Z_0)A(Z_0)u_0,x)
\]
is compact and consists of invertible matrices.  It contains
every $(Z,x)$ with $b_U(Z)\ne0$ and $(Zu,x)\in K$ for some
$u\in N_r$.  This proves the parameter version too.  For
$r=0$ the groups and coordinate spaces are points with their
original unit mass, and $b_U=1$.
\end{proof}

Retain the original normalized level tower.  At the last level
its residual group is $G_0=U(k,1)$, its compact Haar measures
have mass one, and its unipotent measure is the unit mass at
the identity.  These are hypotheses of that prescribed tower,
not new normalizations.  Write the original open chart as
\[
       (a,a',Z,\gamma)\longmapsto
          R(a)^{-1}m(Z,\gamma)R(a'),
       \qquad Z\in\operatorname{GL}_r(\mathbb C),\quad\gamma\in G_0.
\]
Let $P_rF(Z,\gamma)$ denote its prescribed partial Jacquet
average: all right upper roots are integrated, while the
remaining head-lower roots have not yet been integrated.
We use the original radical measure $d\mu_{RC}$ and the
characters $\chi_z(R(a))=\psi_{c_{\rm Std}}(\ell_z(R(a)))$.

\begin{theorem}\label{terminal-transfer-compact-lift}
There are real smooth compact radical bumps $b_L,b_R$ and a
nonzero scalar
\begin{equation}\label{terminal-transfer-kappa}
 \kappa=\left(\int b_L(a)\chi_z(R(a))\,d\mu_{RC}(a)\right)
        \left(\int b_R(a')\overline{\chi_{z'}(R(a'))}
                                              \,d\mu_{RC}(a')\right)
\end{equation}
with the following property.  Suppose $\Phi(Z,\gamma)$ is
smooth on $\det Z\ne0$, $G_0(\gamma)>0$, has the prescribed
left and right compact covariance, and satisfies
\[
       \Phi(Zu,\gamma)=\chi_{z'}(\mu(u))\Phi(Z,\gamma)
                                      \quad(u\in N_r).
\]
If its support modulo this right action lies in a compact
subset of that positive domain, then
\begin{equation}\label{terminal-transfer-chart-lift}
 H\bigl(R(a)^{-1}m(Z,\gamma)R(a')\bigr)
       =b_L(a)b_R(a')b_U(Z)\Phi(Z,\gamma),
\end{equation}
extended by zero, is an original compact coefficient, supported
away from the vanishing of the leading $r$ by $r$ determinant,
and
\begin{equation}\label{terminal-transfer-lift-image}
                   P_rH(Z,\gamma)=\kappa\Phi(Z,\gamma).
\end{equation}
This is an exact identity with the original tower measures.
\end{theorem}

\begin{proof}
Both characters in \eqref{terminal-transfer-kappa} are
continuous and equal to one at zero.  Choose nonnegative
nonzero real bumps supported where their real parts exceed
$1/2$.  Each integral then has positive real part and is
nonzero.  This works simultaneously for the two characters;
no estimate for $\kappa^{-1}$ is used.

The chart and its inverse are smooth on $\det Z\ne0$, by the
already established block Schur-complement formulas.  The
upper-cut lemma makes $b_U\Phi$ compactly supported there,
with positive residual Gram bounded away from zero.
Multiplication by the two radical bumps and extension by
zero therefore give a smooth compact coefficient.  On the
chart the full positive Gram is
$|\det Z|^2G_0(\gamma)$, so it too stays uniformly positive
on this support.

For the integral identity, first perform the right Jacquet
average.  Its exact root-coordinate decomposition consists
of the residual compact group, the residual unipotent group,
$a'$ and $u$.  The residual unipotent group is trivial.
The compact covariance cancels the two representation
matrices, and its normalized Haar mass is one.  The factor
$\chi_{z'}(\mu(u))$ from $\Phi(Zu,\gamma)$ cancels the
inverse Jacquet character.  The $u$ integral is
$\int b_U(Zu)\,du=1$.  The result at
$R(a)^{-1}h^{-1}m(Z,\gamma)$ is consequently
\[
       b_L(a)\kappa_R\,\pi(h^{-1})\Phi(Z,\gamma).
\]
The left compact average cancels $\pi(h^{-1})$ and has mass
one; its remaining radical integral is exactly $\kappa_L$.
Thus the result is \eqref{terminal-transfer-lift-image}.
The radical actions by compact residual isometries preserve
the original real measure, and the ordered root-coordinate
splittings have diagonal Jacobian one.  These are the same
chart and measure decompositions used in the actual
partial-head transfer.  The compact supports above justify
Fubini throughout.
\end{proof}

We next state the operator comparison with all its original
hypotheses.  Fix the place parameters $P$, $1\le l_0\le k$,
$k\ge2$, the normalized tower, the prescribed finite type
projectors $E_\rho$, and the original operator $\Theta$.
Retain both its specified regular radial identity and its
original differential-operator realization.  For every word
$w=((\varepsilon_a,y_a))$ used below, retain the source cut
condition on every $y_a$ and $|y_a-y_b|>2$ for $a\ne b$.
In particular all source denominators at $y_a,y_a\pm1$ and
all shifted pair differences remain invertible.

Write $B_w^{(r)}F$ for the actual tensor family of leading
minor operators at level $r$ and $\mathcal F_w^{(r)}$ for
the prescribed terminal frame operator.  This notation
retains the word order, tensor indices, both signs, every
spectral shift and all the original rational factors.
Let $\mathcal J_rS_h$ be the actual last-level Jacquet
transform of the slice at a head $h$ of size $r$.

\begin{theorem}\label{terminal-transfer-every-head}
For every original compact coefficient $F$, every such word
$w$, every tensor index $e$, every relevant head $h$ of size
$r$, and every $\zeta$ with $|\zeta|^2>TT'$, choose any
last-level $D$-frame $\gamma_\zeta$ with reduced scalar
$\zeta$.  Then
\begin{equation}\label{terminal-transfer-vector}
 V_{\Theta_{w,e}F}(h\mathbin{\|}(1,\zeta))
   =\left[\mathcal F_w^{(r)}
           \bigl(\mathcal J_rS_h(B_w^{(r)}F)\bigr)\right]_e
                                                  (\gamma_\zeta).
\end{equation}
The original scalar normalizations and stabilizer measures
agree exactly.  In particular the original terminal
density equals the original terminal Gram expression,
pointwise at every last-level group element in the positive
domain and for every pair of heads $h,h'$.
\end{theorem}

\begin{proof}
We first construct a single compact coefficient giving the
right-hand side for every head, after a harmless localization
at the chosen head.

Put $\delta_h=\det H_h\ne0$.  Choose a smooth cutoff in
the leading determinant $D_r(g)$, equal to one near
$\delta_h$ and supported away from zero, and call the
resulting coefficient $F_1$.  Every point of the orbit
$h\mathbin{\|}(1,\zeta)$ has $D_r=\delta_h$.
The original $\Theta_{w,e}$ is a differential operator on
the positive domain, so it has the same value there on
$F$ and $F_1$.  The same holds for every component of
$B_w^{(r)}$, including all its derivatives.  Thus the
localization leaves both sides of
\eqref{terminal-transfer-vector} unchanged.  This uses
local equality on an open neighborhood of the orbit,
not merely equality of values on that orbit.

Set $F_e^u=(B_w^{(r)}F_1)_e$ and
$\Psi_e(Z,\gamma)=P_rF_e^u(Z,\gamma)$.
The actual leading-minor operators are differential on
$D_r\ne0$; their coefficients are smooth there.  Locality
therefore keeps each $F_e^u$ in the support of $F_1$, so it
is compact with its leading determinant supported away
from zero.  The chart inverse and compactness show that
$\Psi$ is smooth and compact modulo right $N_r$; it has
exactly the compact covariance and upper character of the
preceding theorem.  More explicitly, a nonzero integrand
puts $(Zu,h^{-1}\gamma h')$ in a fixed compact chart image
of $\operatorname{supp}F_e^u$.  The compact $h,h'$ can be
absorbed into a single compact set.  Conversely, on a
compact positive parameter set, this relation bounds $u$
and both radical coordinates by the chart inverse.  It
justifies smooth dependence and all subsequent integrals.

Apply the terminal operator in the $\gamma$ variable:
\[
       \Phi_e(Z,\gamma)
            =[\mathcal F_w^{(r)}\Psi(Z,\cdot)]_e(\gamma).
\]
This operation is local, preserves the support, and commutes
with integration over a compact head-parameter set.  To
check this for the actual operator, its generators are
multiplication by
\[
 \zeta,\ \zeta^{-1},\ \overline\zeta,
 \ \overline\zeta^{-1},\ G_0,\ G_0^{-1}
\]
and the prescribed finite Lie differential operators on
$G_0$.  All multipliers are smooth on $G_0>0$ because
$TT'>0$ forces $\zeta\ne0$.  Here the finite spectral
algebra is the one already constructed for the actual
terminal Capelli operator.  If $E_g=\langle g^{-1}z,z'\rangle^\perp$
and $D_{ab}=d\pi'_{\mathbb C}(b_a h(\,\cdot\,,b_b))$ for
an orthonormal basis of $E_g$, the exact identities are
\[
 G_0(g)^{-1}(\mathcal T_{22}(Y)\Phi)(g)
       =\Phi(g)Q_{k-1}(D;Y+(k-2)/2),\qquad
 p_\rho(Y)=\prod_{i=1}^{k-1}(Y+\rho_i).
\]
The first identity follows by contracting the two negative
frame vectors in the ordered determinant; right compact
equivariance replaces each remaining Lie derivative by
$D_{ab}$.  Compact branching restricts the surviving fibres
to the common types $\rho$, on which the displayed matrix
acts by $p_\rho(Y)$.  Consequently the original sign
$\sigma=(-1)^{k-l_0}$ gives
\[
 \mathrm{st}.x(Y)=\sigma G_0^{-1}\mathcal T_{22}(Y)
      =\kappa\bigl((\sigma p_\rho(Y))_\rho\bigr),\qquad
 \mathrm{st}.x^{\mathrm{inv}}(Y)
      =\kappa\bigl(((\sigma p_\rho(Y))^{-1})_\rho\bigr).
\]
The second formula is used on the original prescribed
inverse set.  The already constructed $\kappa$ consists
of pointwise smooth compact-type projections and satisfies
these identities on the actual covariant functions
\cite{terminal-transfer-actual-spectrum}.
For a fixed spectral function $f$, interpolation of $t^{-1}$
on its distinct values expresses $\kappa(f^{-1})$ as a
polynomial in $\kappa(f)$.  This covers repeated eigenvalues
and gives precisely the prescribed inverse, rather than
an inverse on a different function space.  Therefore
each entry of $\mathcal F_w^{(r)}$ is a finite sum of
compositions of these local generators.  The induction
over the actual word proves the stated locality and
parameter properties.  No uniform degree or coefficient
bound for this interpolation is asserted or needed.

The union of the finitely many supports modulo $N_r$ is
compact.  Apply the compact lift to $\kappa^{-1}\Phi_e$,
obtaining a compact coefficient $F_2$ with
\begin{equation}\label{terminal-transfer-pullback-F2}
                       P_rF_2=\Phi_e.
\end{equation}
For any relevant head $a$ of size $r$, the exact partial
unfolding is
\[
 \mathcal J_rS_aG(\gamma)
      =\int_{L_a}\psi_{\beta_U}(\theta_-(L))
                       P_rG(LH_a,\gamma)\,dL,
\]
where $L_a$ consists of the original restricted head-lower
roots.  This formula includes the nonsingleton blocks of
$a$.  Compact support modulo right $N_r$ makes this lower
integral compactly supported: if $LH_a u$ belongs to the
compact support representative, block Gaussian elimination
bounds $L,u$.  The actual root quotient removes all internal
stabilizer directions.  The support bound is also uniform
for $\gamma$ in a compact positive neighborhood, so the
finite terminal differential operators pass through the
integral.  Equation \eqref{terminal-transfer-pullback-F2}
therefore gives, for every such $a$,
\begin{equation}\label{terminal-transfer-all-slices}
 \mathcal J_rS_aF_2
       =\left[\mathcal F_w^{(r)}
                 \bigl(\mathcal J_rS_a(B_w^{(r)}F_1)\bigr)
                                                \right]_e.
\end{equation}

On the regular heads this is exactly the regular radial
picture prescribed for $\Theta_{w,e}F_1$.  We spell out
the compatibility: restriction of the last group to its
embedded $U(k,1)$ preserves $\zeta,G_0,T,T'$ and the six
displayed multiplier functions, and restriction intertwines
every terminal Lie direction.  Hence it intertwines the
word operator, including its polynomially realized inverse
spectral factors.  The leading-minor part contributes the
same previously proved radial gauge on both sides,
\[
                  \prod_{j=0}^{r-1}|c_j|^{2(r-j)+k}.
\]
The exact leading-minor radial identity thus identifies the
regular radial transform of $F_2$ with that of
$\Theta_{w,e}F_1$, with no scalar left over.

The differential realization of $\Theta_{w,e}F_1$ has a
compact smooth representative on the full unitary group.
Indeed its coefficients are smooth on $G>0$, all derivatives
of $F_1$ vanish outside its compact support, and that support
lies in $G\ge\epsilon$.  Extend the differential expression
by zero across $G=0$.  This compact representative agrees
with the original operator on the positive domain and
hence in every nonzero orbital vector.  Subtract $F_2$.
The difference has zero regular radial transform, so the
every-head density corollary of
Section~\ref{actual-unitary-density} gives
\[
             V_{\Theta_{w,e}F_1}(b)=V_{F_2}(b)
                                  \quad\hbox{for every head }b.
\]
Finally the exact root-coordinate factorization at an
appended singleton gives
\[
      V_{F_2}(h\mathbin{\|}(1,\zeta))
            =\mathcal J_rS_hF_2(\gamma_\zeta).
\]
At rank zero the unipotent group is the identity, so this
last assertion is the same normalized compact average.
At positive rank it is the same radical and restricted
root disintegration already used above.  All base frames
are the prescribed embedded frames.  Combining with
\eqref{terminal-transfer-all-slices} and undoing the
localization proves \eqref{terminal-transfer-vector}.

For completeness we verify the terminal weights.  At
$b=h\mathbin{\|}(1,\zeta)$ the source weight is
\[
 w_{\rm pt}(b)=\frac{|\det H_b|^2}{G(H_b)}
                  =\frac{|\zeta|^2}{|\zeta|^2-TT'},
\]
independent of $h$.  For type $\rho$ and tensor index $e$
let $c_{w,e,\rho}$ be the original positive rational word
product, with its unchanged plus and minus recursion.
The diagonal terminal operator acts on that type by
\[
       \lambda_{w,e,\rho}(\zeta)
         =\left(\frac{|\zeta|^2}{|\zeta|^2-TT'}\right)^{|w|}
                    c_{w,e,\rho}=d_H(b,e,\rho)^2.
\]
This follows by diagonalizing the finite compact type
algebra: select a real spectral parameter separating its
finitely many distinct type polynomials, then interpolate
the prescribed word eigenvalues there.  The original
interlacing and cut conditions give
$c_{w,e,\rho}>0$, so the square root in $d_H$ is the
positive one.  The orbital vector is fixed by the original
diagonal compact stabilizer of the remaining positive
complement, by its defining compact average.  The prescribed
compact branching decomposition and type-projector
completeness therefore apply on this vector; completeness
is not assumed on an unrelated larger representation space.
These type projections are mutually orthogonal and
self-adjoint.  Thus
the terminal density is precisely
\[
 \sum_{e,\rho}\lambda_{w,e,\rho}(\zeta)
   \left\langle E_\rho V_{\Theta_{w,e}F}(h\mathbin{\|}(1,\zeta)),
       E_\rho V_{\Theta_{w,e}F}(h'\mathbin{\|}(1,\zeta))
   \right\rangle_{\rm HS},
\]
which is the original terminal Gram expression, since
$d_H(h\mathbin{\|}(1,\zeta),e,\rho)
 =d_H(h'\mathbin{\|}(1,\zeta),e,\rho)$.
For the empty word the operators and weights are identities
and the expression is the original Hilbert--Schmidt pairing.
Compact covariance and unitary invariance of this pairing
extend the identity from the chosen $D$-frame to every
last-level group element with the same $\zeta$.
\end{proof}

\paragraph{Closed interface and remaining input.}
This proves the original compact top-lift and terminal
transfer contracts, including all block heads, all admissible
words, the empty word, and the rank-zero case
\cite{terminal-transfer-source}.  The existence of the
original place and operator with its prescribed regular
radial identity and uniform differential bounds remains
exactly a hypothesis where the source places it.  No estimate
for the auxiliary cutoffs, interpolation degree or
$\kappa^{-1}$ has been inserted into an arithmetic bound;
they cancel in exact identities.  The terminal transfer
alone does not establish the global weighted kernel estimate
or complete W06/W08.
